\documentclass[../../main.tex]{subfiles}% 注意这里的文件路径不能用 ./main.tex ,否则用latexmk编译子文件会报错
\graphicspath{{\subfix{./image/}}{\subfix{../../image/}}} % 指定图片引用路径,后续可以直接使用图片文件名
% 如果图片存放在上级目录,则使用 ../../image/ 作为备用路径

% 例如:
% \begin{figure}[H]
% \centering
% \includegraphics[width=0.3\textwidth]{图.png}
% \caption{}
% \label{figure:图}
% \end{figure}
% 注意:上述\label{}一定要放在\caption{}之后,否则引用图片序号会只会显示??.

\begin{document}
% SOURCE p=85; printed=70; object=ch02-nd-lebesgue-title; region=top
\section{\texorpdfstring{$n$}{n} 维 Lebesgue 积分}\label{section:ch02-n-dimensional-lebesgue-integral}
% SOURCE p=85; printed=70; object=ch02-nd-lebesgue-measure; region=upper
\begin{definition}[$n$ 维 Lebesgue 测度]\label{definition:ch02-nd-lebesgue-measure}
$\mathbb{R}^n$ 上的 \textbf{Lebesgue 测度} $m^n$ 是 $\mathbb{R}$ 上的 Lebesgue 测度与自身的 $n$ 重乘积的完备化，即 $\mathscr{B}_{\mathbb{R}}\otimes\cdots\otimes\mathscr{B}_{\mathbb{R}}=\mathscr{B}_{\mathbb{R}^n}$ 上的 $m\times\cdots\times m$ 的完备化；等价地，它是 $\mathscr{L}\otimes\cdots\otimes\mathscr{L}$ 上的 $m\times\cdots\times m$ 的完备化。$m^n$ 的定义域 $\mathscr{L}^n$ 是 $\mathbb{R}^n$ 中的 \textbf{Lebesgue 可测集}组成的类。
\end{definition}
% SOURCE p=85; printed=70; object=ch02-nd-lebesgue-notation; region=upper
\begin{remark}
有时也将 $m^n$ 看作较小定义域 $\mathscr{B}_{\mathbb{R}^n}$ 上的测度。不致混淆时，通常省略上标 $n$，以 $m$ 表示 $m^n$；与 $n=1$ 时一样，通常用 $\int f(x)\,\mathrm{d}x$ 表示 $\int f\,\mathrm{d}m$。
\end{remark}
% SOURCE p=85; printed=70; object=ch02-nd-regularity-intro; region=upper-middle
首先将\cref{section:ch01-borel-measures-real-line}中的一些结果推广到 $n$ 维情形。
% SOURCE p=85; printed=70; object=ch02-rectangle-sides; region=upper-middle
\begin{definition}[矩形的边]\label{definition:ch02-rectangle-sides}
以下若 $E=\prod_{j=1}^n E_j$ 为 $\mathbb{R}^n$ 中的矩形，就将集合 $E_j\subseteq\mathbb{R}$ 称为 $E$ 的\textbf{边}。
\end{definition}
% SOURCE p=85; printed=70; object=ch02-nd-regularity; number=2.40; region=middle
\begin{theorem}\label{theorem:ch02-nd-regularity}
设 $E\in\mathscr{L}^n$。
\begin{enumerate}
\item\label{item:ch02-nd-regularity-open-compact} $m(E)=\inf\{m(U):U\supseteq E,\ U\text{ 为开集}\}=\sup\{m(K):K\subseteq E,\ K\text{ 为紧集}\}$。
\item\label{item:ch02-nd-regularity-borel-null} $E=A_1\cup N_1=A_2\setminus N_2$，其中 $A_1$ 为 $F_\sigma$ 集，$A_2$ 为 $G_\delta$ 集，且 $m(N_1)=m(N_2)=0$。
\item\label{item:ch02-nd-regularity-rectangles} 若 $m(E)<\infty$，则对每个 $\epsilon>0$，存在有限个两两不交、各边均为区间的矩形 $\{R_j\}_1^N$，使 $m(E\mathbin{\triangle}\bigcup_1^N R_j)<\epsilon$。
\end{enumerate}
\end{theorem}
% SOURCE p=85; printed=70; object=ch02-proof-nd-regularity; region=middle-lower
\begin{proof}
由乘积测度的定义，若 $E\in\mathscr{L}^n$ 且 $\epsilon>0$，则存在可数个矩形 $\{T_j\}$，使 $E\subseteq\bigcup_1^\infty T_j$，且 $\sum_1^\infty m(T_j)\le m(E)+\epsilon$。对每个 $j$，将\cref{theorem:ch01-stieltjes-regularity}应用于 $T_j$ 的各边，可以找到各边为开集的矩形 $U_j\supseteq T_j$，使 $m(U_j)<m(T_j)+\epsilon2^{-j}$。若 $U=\bigcup_1^\infty U_j$，则 $U$ 为开集，且 $m(U)\le\sum_1^\infty m(U_j)<m(E)+2\epsilon$。这就证明了\subcref{theorem:ch02-nd-regularity}{item:ch02-nd-regularity-open-compact}中的第一个等式；第二个等式和\subcref{theorem:ch02-nd-regularity}{item:ch02-nd-regularity-borel-null}可仿照\cref{theorem:ch01-stieltjes-regularity,theorem:ch01-stieltjes-measurable-sets}的证明得到。其次，若 $m(E)<\infty$，则对所有 $j$，有 $m(U_j)<\infty$。由于 $U_j$ 的各边是开区间的可数并，取适当的有限子并，可得到各边为有限个区间的并的矩形 $V_j\subseteq U_j$，使 $m(V_j)\ge m(U_j)-\epsilon2^{-j}$。当 $N$ 充分大时，
\begin{align*}
m\left(E\setminus\bigcup_1^N V_j\right)
\le m\left(\bigcup_1^N(U_j\setminus V_j)\right)
+m\left(\bigcup_{N+1}^\infty U_j\right)<2\epsilon,
\end{align*}
并且
\begin{align*}
m\left(\bigcup_1^N V_j\setminus E\right)<m\left(\bigcup_1^\infty U_j\setminus E\right)<\epsilon,
\end{align*}
从而 $m(E\mathbin{\triangle}\bigcup_1^N V_j)<3\epsilon$。由于 $\bigcup_1^N V_j$ 可以表示为有限个两两不交、各边为区间的矩形的并，\subcref{theorem:ch02-nd-regularity}{item:ch02-nd-regularity-rectangles}得证。
\end{proof}
% SOURCE p=85; printed=70; object=ch02-nd-lone-density; number=2.41; region=bottom
\begin{theorem}\label{theorem:ch02-nd-lone-density}
若 $f\in L^1(m)$ 且 $\epsilon>0$，则存在简单函数 $\phi=\sum_1^N a_j\chi_{R_j}$，其中每个 $R_j$ 为区间的乘积，使 $\int|f-\phi|<\epsilon$；并且存在在某个有界集之外为零的连续函数 $g$，使 $\int|f-g|<\epsilon$。
\end{theorem}
% SOURCE p=86; printed=71; object=ch02-proof-nd-lone-density; region=top
\begin{proof}
仿照\cref{theorem:ch02-lone-density}的证明，用简单函数逼近 $f$，再用\subcref{theorem:ch02-nd-regularity}{item:ch02-nd-regularity-rectangles}，以所需形式的函数逼近这些简单函数。最后，对\cref{theorem:ch02-lone-density}证明中的论证作显然的推广，就能以连续函数逼近这样的 $\phi$。
\end{proof}
% SOURCE p=86; printed=71; object=ch02-nd-translation; number=2.42; region=upper
\begin{theorem}[Lebesgue 测度的平移不变性]\label{theorem:ch02-nd-translation}
Lebesgue 测度具有平移不变性。更确切地说，\begin{definition*}[平移映射]
对 $a\in\mathbb{R}^n$，定义 $\tau_a:\mathbb{R}^n\to\mathbb{R}^n$ 为 $\tau_a(x)=x+a$。
\end{definition*}
\begin{enumerate}
\item\label{item:ch02-nd-translation-sets} 若 $E\in\mathscr{L}^n$，则 $\tau_a(E)\in\mathscr{L}^n$，且 $m(\tau_a(E))=m(E)$。
\item\label{item:ch02-nd-translation-functions} 若 $f:\mathbb{R}^n\to\mathbb{C}$ 为 Lebesgue 可测函数，则 $f\circ\tau_a$ 也为 Lebesgue 可测函数。此外，若 $f\ge0$ 或 $f\in L^1(m)$，则 $\int(f\circ\tau_a)\,\mathrm{d}m=\int f\,\mathrm{d}m$。
\end{enumerate}
\end{theorem}
% SOURCE p=86; printed=71; object=ch02-proof-nd-translation; region=middle
\begin{proof}
由于 $\tau_a$ 及其逆映射 $\tau_{-a}$ 连续，它们保持 Borel 集类。若 $E$ 为矩形，公式 $m(\tau_a(E))=m(E)$ 容易由一维结果（\cref{theorem:ch01-lebesgue-invariance}）得到；由于 $m$ 由其在矩形上的作用唯一确定（\cref{theorem:ch01-measure-extension}中的唯一性），这个公式随即对一般 Borel 集成立。特别地，满足 $m(E)=0$ 的 Borel 集 $E$ 的全体在 $\tau_a$ 下不变。\subcref{theorem:ch02-nd-translation}{item:ch02-nd-translation-sets}现在立即成立。

若 $f$ 为 Lebesgue 可测函数，$B$ 为 $\mathbb{C}$ 中的 Borel 集，则 $f^{-1}(B)=E\cup N$，其中 $E$ 为 Borel 集，$m(N)=0$。但 $\tau_a^{-1}(E)$ 为 Borel 集，且 $m(\tau_a^{-1}(N))=0$，故 $(f\circ\tau_a)^{-1}(B)\in\mathscr{L}^n$，从而 $f\circ\tau_a$ 为 Lebesgue 可测函数。当 $f=\chi_E$ 时，等式 $\int(f\circ\tau_a)\,\mathrm{d}\mu=\int f\,\mathrm{d}\mu$ 归结为 $m(\tau_{-a}(E))=m(E)$。由线性性，它对简单函数成立；再由积分定义，它对非负可测函数成立。最后，取实部和虚部的正部、负部，就得到 $f\in L^1(m)$ 的结论。
\end{proof}
% SOURCE p=86; printed=71; object=ch02-jordan-intro; region=middle-lower
现在将 $\mathbb{R}^n$ 上的 Lebesgue 测度与高等微积分教材中通常介绍的较朴素的 $n$ 维测度理论作比较。
% SOURCE p=86; printed=71; object=ch02-cube; region=middle-lower
\begin{definition}[立方体]\label{definition:ch02-cube}
在这里，$\mathbb{R}^n$ 中的\textbf{立方体}是 $n$ 个闭区间的笛卡尔积，且这些区间的长度均相等。
\end{definition}
% SOURCE p=86; printed=71; object=ch02-dyadic-cubes; region=lower
\begin{definition}[二进立方体族]\label{definition:ch02-dyadic-cubes}
对 $k\in\mathbb{Z}$，令 $\mathscr{Q}_k$ 为边长为 $2^{-k}$、顶点在格点集 $(2^{-k}\mathbb{Z})^n$ 中的立方体组成的族。也就是说，$\prod_{j=1}^n[a_j,b_j]\in\mathscr{Q}_k$ 当且仅当对所有 $j$，$2^ka_j$ 和 $2^kb_j$ 为整数，且 $b_j-a_j=2^{-k}$。
\end{definition}
% SOURCE p=86; printed=71; object=ch02-dyadic-cube-properties; region=lower
\begin{remark}
$\mathscr{Q}_k$ 中任意两个不同立方体的内部互不相交；将 $\mathscr{Q}_k$ 中各立方体的边二等分，就得到 $\mathscr{Q}_{k+1}$ 中的立方体。
\end{remark}
% SOURCE p=86; printed=71; object=ch02-cube-approximations; region=lower
\begin{definition}[立方体网格的内、外逼近]\label{definition:ch02-cube-approximations}
若 $E\subseteq\mathbb{R}^n$，定义立方体网格 $\mathscr{Q}_k$ 对 $E$ 的\textbf{内逼近}和\textbf{外逼近}为
\begin{align*}
\underline{A}(E,k)&=\bigcup\{Q\in\mathscr{Q}_k:Q\subseteq E\},\\
\overline{A}(E,k)&=\bigcup\{Q\in\mathscr{Q}_k:Q\cap E\ne\varnothing\}.
\end{align*}
见\cref{figure:ch02-jordan-approximation}。
\end{definition}
% SOURCE p=87; printed=72; object=ch02-jordan-approximation-figure; number=2.2; region=top
\begin{figure}[H]
\centering
\includegraphics[width=0.48\textwidth]{ch02-jordan-approximation.png}
\caption{集合的内、外含量的逼近。}\label{figure:ch02-jordan-approximation}
\end{figure}
% SOURCE p=86; printed=71; object=ch02-cube-approximation-measures; region=bottom
\begin{remark}
$\underline{A}(E,k)$ 的测度（无论在朴素几何意义下，还是在 Lebesgue 意义下）就是 $2^{-nk}$ 乘以包含于 $\underline{A}(E,k)$ 的 $\mathscr{Q}_k$ 中的立方体个数，将它记作 $m(\underline{A}(E,k))$；$m(\overline{A}(E,k))$ 的情形相同。此外，集合 $\underline{A}(E,k)$ 随 $k$ 递增，而 $\overline{A}(E,k)$ 随 $k$ 递减，因为 $\mathscr{Q}_k$ 中每个立方体都是 $\mathscr{Q}_{k+1}$ 中一些立方体的并。
\end{remark}
% SOURCE p=86,87; printed=71,72; object=ch02-jordan-content; region=bottom/middle
\begin{definition}[内含量、外含量与 Jordan 含量]\label{definition:ch02-jordan-content}
极限
\begin{align*}
\underline{\kappa}(E)=\lim_{k\to\infty}m(\underline{A}(E,k)),\qquad
\overline{\kappa}(E)=\lim_{k\to\infty}m(\overline{A}(E,k))
\end{align*}
均存在，分别称为 $E$ 的\textbf{内含量}和\textbf{外含量}。若二者相等，其共同值 $\kappa(E)$ 就是 $E$ 的 \textbf{Jordan 含量}。
\end{definition}
% SOURCE p=87; printed=72; object=ch02-jordan-content-remarks; region=middle
\begin{remark}
有两点说明。第一，通常用各边为区间的一般矩形来定义 Jordan 含量，而不采用这里的二进立方体，但结果相同。第二，虽然上述定义对任意 $E\subseteq\mathbb{R}^n$ 都有意义，Jordan 含量理论只有在 $E$ 有界时才有实质意义，否则 $\overline{\kappa}(E)$ 总等于 $\infty$。
\end{remark}
% SOURCE p=87; printed=72; object=ch02-jordan-limit-sets; region=middle
令
\begin{align*}
\underline{A}(E)=\bigcup_1^\infty\underline{A}(E,k),\qquad
\overline{A}(E)=\bigcap_1^\infty\overline{A}(E,k).
\end{align*}
% SOURCE p=87; printed=72; object=ch02-jordan-lebesgue-criterion; region=middle-lower
\begin{proposition}\label{proposition:ch02-jordan-lebesgue-criterion}
有 $\underline{A}(E)\subseteq E\subseteq\overline{A}(E)$，$\underline{A}(E)$ 和 $\overline{A}(E)$ 均为 Borel 集，且 $\underline{\kappa}(E)=m(\underline{A}(E))$、$\overline{\kappa}(E)=m(\overline{A}(E))$。因此，$E$ 的 Jordan 含量存在，当且仅当 $m(\overline{A}(E)\setminus\underline{A}(E))=0$；这蕴含 $E$ 为 Lebesgue 可测集，且 $m(E)=\kappa(E)$。
\end{proposition}
% SOURCE p=87; printed=72; object=ch02-open-cubes-intro; region=lower
为进一步说明 Lebesgue 测度与产生 Jordan 含量的逼近过程之间的关系，我们证明下面的引理。（引理的第二部分以后还会用到。）
% SOURCE p=87; printed=72; object=ch02-open-cubes; number=2.43; region=lower
\begin{lemma}\label{lemma:ch02-open-cubes}
若 $U\subseteq\mathbb{R}^n$ 为开集，则 $U=\underline{A}(U)$。此外，$U$ 是可数个内部两两不交的立方体的并。
\end{lemma}
% SOURCE p=87; printed=72; object=ch02-proof-open-cubes; region=lower-bottom
\begin{proof}
若 $x\in U$，令 $\delta=\inf\{|y-x|:y\notin U\}$；由于 $U$ 开，$\delta>0$。若 $Q$ 为 $\mathscr{Q}_k$ 中包含 $x$ 的一个立方体，则每个 $y\in Q$ 到 $x$ 的距离至多为 $2^{-k}\sqrt n$；最坏情形是对所有 $j$，$|x_j-y_j|=2^{-k}$。因此，只要 $k$ 大到使 $2^{-k}\sqrt n<\delta$，就有 $Q\subseteq U$。于是 $x\in\underline{A}(U,k)\subseteq\underline{A}(U)$。

这说明 $\underline{A}(U)=U$。第二个断言可通过写成 $\underline{A}(U)=\underline{A}(U,0)\cup\bigcup_1^\infty[\underline{A}(U,k)\setminus\underline{A}(U,k-1)]$ 得到。$\underline{A}(U,0)$ 是 $\mathscr{Q}_0$ 中一些立方体的（可数）并；而当 $k\ge1$ 时，$\underline{A}(U,k)\setminus\underline{A}(U,k-1)$ 的闭包是 $\mathscr{Q}_k$ 中一些立方体的（可数）并。这些立方体的内部均两两不交，结论即得。
\end{proof}
% SOURCE p=88; printed=73; object=ch02-open-inner-content; region=top
\begin{proposition}\label{proposition:ch02-open-inner-content}
任意开集的 Lebesgue 测度等于其内含量。
\end{proposition}
\begin{proof}
这由\cref{lemma:ch02-open-cubes}立即得出。
\end{proof}
% SOURCE p=88; printed=73; object=ch02-compact-outer-content; region=upper
\begin{proposition}\label{proposition:ch02-compact-outer-content}
任意紧集的 Lebesgue 测度等于其外含量。
\end{proposition}
\begin{proof}
设 $F\subseteq\mathbb{R}^n$ 为紧集。可取一个大立方体，例如 $Q_0=\{x:\max|x_j|\le2^M\}$，使其内部 $\operatorname{int}(Q_0)$ 包含 $F$。若 $Q\in\mathscr{Q}_k$ 且 $Q\subseteq Q_0$，则或者 $Q\cap F\ne\varnothing$，或者 $Q\subseteq Q_0\setminus F$，因而 $m(\overline{A}(F,k))+m(\underline{A}(Q_0\setminus F,k))=m(Q_0)$。令 $k\to\infty$，得到 $\overline{\kappa}(F)+\underline{\kappa}(Q_0\setminus F)=m(Q_0)$。但 $Q_0\setminus F$ 是开集 $\operatorname{int}(Q_0)\setminus F$ 与 $Q_0$ 的边界之并；后者含量为零，所以 $\underline{\kappa}(Q_0\setminus F)=\underline{\kappa}(\operatorname{int}(Q_0)\setminus F)=m(Q_0\setminus F)$。结论随即成立。
\end{proof}
% SOURCE p=88; printed=73; object=ch02-jordan-lebesgue-comparison; region=upper-middle
\begin{remark}
结合这些结果与\subcref{theorem:ch02-nd-regularity}{item:ch02-nd-regularity-open-compact}，可以准确看出 Lebesgue 测度与 Jordan 含量的关系。$E$ 的 Jordan 含量通过有限个立方体的并从内部和外部逼近 $E$ 来定义。另一方面，$E$ 的 Lebesgue 测度由两步逼近过程给出：先用开集从外部、用紧集从内部逼近 $E$；然后用有限个立方体的并从内部逼近这些开集、从外部逼近这些紧集。Lebesgue 可测集恰好是使这种“先外后内”和“先内后外”的逼近在极限中给出同一结果的集合。（比较\cref{section:ch01-outer-measures}中的\cref{exercise:ch01-inner-measure}。）
\end{remark}
% SOURCE p=88; printed=73; object=ch02-linear-change-intro; region=middle
现在考察 Lebesgue 积分在线性变换下的行为。
% SOURCE p=88; printed=73; object=ch02-linear-map-matrix; region=middle
\begin{definition}[线性映射的矩阵与行列式]\label{definition:ch02-linear-map-matrix}
将线性映射 $T:\mathbb{R}^n\to\mathbb{R}^n$ 与矩阵 $(T_{ij})=(e_i\cdot Te_j)$ 等同，其中 $\{e_j\}$ 为 $\mathbb{R}^n$ 的标准基。这个矩阵的行列式记作 $\det T$。
\end{definition}
% SOURCE p=88; printed=73; object=ch02-determinant-composition; region=middle
\begin{remark}
回顾 $\det(T\circ S)=(\det T)(\det S)$。
\end{remark}
% SOURCE p=88; printed=73; object=ch02-general-linear-group; region=middle
\begin{definition}[一般线性群]\label{definition:ch02-general-linear-group}
用标准记号 $GL(n,\mathbb{R})$ 表示 $\mathbb{R}^n$ 上可逆线性变换组成的群，称为\textbf{一般线性群}。
\end{definition}
% SOURCE p=88; printed=73; object=ch02-elementary-linear-transformations; region=middle-lower
\begin{definition}[三类初等线性变换]\label{definition:ch02-elementary-linear-transformations}
这里使用三类“初等”变换：第一类将一个坐标乘以非零常数 $c$，其余坐标不变；第二类将一个坐标的倍数加到另一个坐标上，除后一个坐标外其余坐标不变；第三类交换两个坐标，其余坐标不变。用符号表示为
\begin{align*}
T_1(x_1,\ldots,x_j,\ldots,x_n)&=(x_1,\ldots,cx_j,\ldots,x_n)\quad(c\ne0),\\
T_2(x_1,\ldots,x_j,\ldots,x_n)&=(x_1,\ldots,x_j+cx_k,\ldots,x_n)\quad(k\ne j),\\
T_3(x_1,\ldots,x_j,\ldots,x_k,\ldots,x_n)&=(x_1,\ldots,x_k,\ldots,x_j,\ldots,x_n).
\end{align*}
\end{definition}
% SOURCE p=88; printed=73; object=ch02-elementary-factorization; region=lower
\begin{proposition}\label{proposition:ch02-elementary-factorization}
每个 $T\in GL(n,\mathbb{R})$ 都可写为有限个上述三类初等变换的乘积。
\end{proposition}
\begin{proof}
这就是每个非奇异矩阵都能通过行变换化为单位矩阵这一事实。
\end{proof}
% SOURCE p=88; printed=73; object=ch02-linear-change; number=2.44; region=lower
\begin{theorem}\label{theorem:ch02-linear-change}
设 $T\in GL(n,\mathbb{R})$。
\begin{enumerate}
\item\label{item:ch02-linear-change-integrals} 若 $f$ 为 $\mathbb{R}^n$ 上的 Lebesgue 可测函数，则 $f\circ T$ 也为 Lebesgue 可测函数。若 $f\ge0$ 或 $f\in L^1(m)$，则
% SOURCE p=88; printed=73; object=ch02-linear-change-equation; number=2.45; region=bottom
\begin{align}\label{align:ch02-linear-change}
\int f(x)\,\mathrm{d}x=|\det T|\int f\circ T(x)\,\mathrm{d}x.
\end{align}
% SOURCE p=88; printed=73; object=ch02-linear-change-sets; region=bottom
\item\label{item:ch02-linear-change-sets} 若 $E\in\mathscr{L}^n$，则 $T(E)\in\mathscr{L}^n$，且 $m(T(E))=|\det T|m(E)$。
\end{enumerate}
\end{theorem}
% SOURCE p=89; printed=74; object=ch02-proof-linear-change; region=upper
\begin{proof}
先假设 $f$ 为 Borel 可测函数。由于 $T$ 连续，$f\circ T$ 也为 Borel 可测函数。若\eqref{align:ch02-linear-change}对变换 $T$ 和 $S$ 成立，它也对 $T\circ S$ 成立，因为
\begin{align*}
\int f(x)\,\mathrm{d}x
&=|\det T|\int f\circ T(x)\,\mathrm{d}x\\
&=|\det T||\det S|\int(f\circ T)\circ S(x)\,\mathrm{d}x\\
&=|\det(T\circ S)|\int f\circ(T\circ S)(x)\,\mathrm{d}x.
\end{align*}
因而，只需在 $T$ 属于上述 $T_1,T_2,T_3$ 三类时证明\eqref{align:ch02-linear-change}。但这只是 Fubini--Tonelli 定理的一个简单推论。对于 $T_3$，交换变量 $x_j$ 和 $x_k$ 的积分次序；对于 $T_1$ 和 $T_2$，先关于 $x_j$ 积分，并利用一维公式
\begin{align*}
\int f(t)\,\mathrm{d}t=|c|\int f(ct)\,\mathrm{d}t,\qquad
\int f(t+a)\,\mathrm{d}t=\int f(t)\,\mathrm{d}t,
\end{align*}
它们由\cref{theorem:ch01-lebesgue-invariance}得到。容易验证 $\det T_1=c$、$\det T_2=1$、$\det T_3=-1$，所以\eqref{align:ch02-linear-change}得证。此外，若 $E$ 为 Borel 集，则 $T(E)$ 也为 Borel 集，因为 $T^{-1}$ 连续；取 $f=\chi_{T(E)}$，得到 $m(T(E))=|\det T|m(E)$。特别地，Borel 零测集类在 $T$ 和 $T^{-1}$ 下不变，从而 $\mathscr{L}^n$ 也不变。Lebesgue 可测函数和集合的结论现在可仿照\cref{theorem:ch02-nd-translation}的证明得到。
\end{proof}
% SOURCE p=89; printed=74; object=ch02-rotation-invariance; number=2.46; region=middle
\begin{corollary}\label{corollary:ch02-rotation-invariance}
Lebesgue 测度在旋转下不变。
\end{corollary}
\begin{proof}
旋转是满足 $TT^*=I$ 的线性映射，其中 $T^*$ 为 $T$ 的转置。由于 $\det T=\det T^*$，这个条件蕴含 $|\det T|=1$。
\end{proof}
% SOURCE p=89; printed=74; object=ch02-differentiable-change-intro; region=middle
下面将\cref{theorem:ch02-linear-change}推广到可微映射。这一结果不会在本书其他地方使用，初读时可以略去。在\cref{section:ch11-change-of-variables}中，我们将用稍有不同的方法证明其一个推广。
% SOURCE p=89; printed=74; object=ch02-derivative-map; region=middle-lower
\begin{definition}[微分映射]\label{definition:ch02-derivative-map}
设 $G=(g_1,\ldots,g_n)$ 是从开集 $\Omega\subseteq\mathbb{R}^n$ 到 $\mathbb{R}^n$ 的映射，其各分量 $g_j$ 属于 $C^1$ 类，即具有连续的一阶偏导数。用 $D_xG$ 表示由点 $x$ 处偏导数矩阵 $((\partial g_i/\partial x_j)(x))$ 确定的线性映射。
\end{definition}
% SOURCE p=89; printed=74; object=ch02-linear-derivative-remark; region=lower
\begin{remark}
若 $G$ 为线性映射，则对所有 $x$，$D_xG=G$。
\end{remark}
% SOURCE p=89; printed=74; object=ch02-c-one-diffeomorphism; region=lower
\begin{definition}[$C^1$ 微分同胚]\label{definition:ch02-c-one-diffeomorphism}
若 $G$ 为单射，且对所有 $x\in\Omega$，$D_xG$ 可逆，则称 $G$ 为 \textbf{$C^1$ 微分同胚}。
\end{definition}
% SOURCE p=89; printed=74; object=ch02-inverse-diffeomorphism; region=lower
\begin{proposition}\label{proposition:ch02-inverse-diffeomorphism}
此时，$G^{-1}:G(\Omega)\to\Omega$ 也为 $C^1$ 微分同胚，并且对所有 $x\in G(\Omega)$，有 $D_x(G^{-1})=[D_{G^{-1}(x)}G]^{-1}$。
\end{proposition}
\begin{proof}
这是逆函数定理的结论。
\end{proof}
% SOURCE p=89; printed=74; object=ch02-differentiable-change; number=2.47; region=lower-bottom
\begin{theorem}\label{theorem:ch02-differentiable-change}
设 $\Omega$ 为 $\mathbb{R}^n$ 中的开集，$G:\Omega\to\mathbb{R}^n$ 为 $C^1$ 微分同胚。
\begin{enumerate}
\item\label{item:ch02-differentiable-change-integrals} 若 $f$ 为 $G(\Omega)$ 上的 Lebesgue 可测函数，则 $f\circ G$ 为 $\Omega$ 上的 Lebesgue 可测函数。若 $f\ge0$ 或 $f\in L^1(G(\Omega),m)$，则
\begin{align*}
\int_{G(\Omega)}f(x)\,\mathrm{d}x=\int_\Omega f\circ G(x)|\det D_xG|\,\mathrm{d}x.
\end{align*}
\item\label{item:ch02-differentiable-change-sets} 若 $E\subseteq\Omega$ 且 $E\in\mathscr{L}^n$，则 $G(E)\in\mathscr{L}^n$，且 $m(G(E))=\int_E|\det D_xG|\,\mathrm{d}x$。
\end{enumerate}
\end{theorem}
% SOURCE p=90; printed=75; object=ch02-proof-differentiable-change; region=top
\begin{proof}
只需考虑 Borel 可测函数和 Borel 集。由于 $G$ 和 $G^{-1}$ 都连续，这种情形下不存在可测性问题；一般情形则仿照\cref{theorem:ch02-nd-translation}的证明得到。
\begin{definition*}[本证明中的范数记号]
对 $x\in\mathbb{R}^n$ 和 $T=(T_{ij})\in GL(n,\mathbb{R})$，令
\begin{align*}
\|x\|=\max_{1\le j\le n}|x_j|,\qquad
\|T\|=\max_{1\le i\le n}\sum_{j=1}^n|T_{ij}|.
\end{align*}
\end{definition*}
于是 $\|Tx\|\le\|T\|\|x\|$，且 $\{x:\|x-a\|\le h\}$ 是以 $a$ 为中心、边长为 $2h$ 的立方体。

设 $Q$ 为 $\Omega$ 中的立方体，例如 $Q=\{x:\|x-a\|\le h\}$。由中值定理，对连接 $x$ 和 $a$ 的线段上的某个 $y$，有 $g_j(x)-g_j(a)=\sum_j(x_j-a_j)(\partial g_j/\partial x_j)(y)$，从而对 $x\in Q$，有 $\|G(x)-G(a)\|\le h(\sup_{y\in Q}\|D_yG\|)$。换言之，$G(Q)$ 包含于一个立方体中，后者边长为 $Q$ 的边长的 $\sup_{y\in Q}\|D_yG\|$ 倍。因此，由\cref{theorem:ch02-linear-change}，$m(G(Q))\le(\sup_{y\in Q}\|D_yG\|)^n m(Q)$。若 $T\in GL(n,\mathbb{R})$，以 $T^{-1}\circ G$ 代替 $G$ 应用这个公式，并结合\cref{theorem:ch02-linear-change}，得到
% SOURCE p=90; printed=75; object=ch02-cube-image-estimate; number=2.48; region=middle
\begin{align}\label{align:ch02-cube-image-estimate}
m(G(Q))&=|\det T|m(T^{-1}(G(Q)))\notag\\
&\le|\det T|\left(\sup_{y\in Q}\|T^{-1}D_yG\|\right)^n m(Q).
\end{align}
% SOURCE p=90,91; printed=75,76; object=ch02-proof-differentiable-change-continued; region=middle-bottom/top
由于 $D_yG$ 关于 $y$ 连续，对任意 $\epsilon>0$，可选取 $\delta>0$，使得若 $y,z\in Q$ 且 $\|y-z\|\le\delta$，则 $\|(D_zG)^{-1}D_yG\|^n\le1+\epsilon$。现在将 $Q$ 分成子立方体 $Q_1,\ldots,Q_N$，其内部两两不交，边长至多为 $\delta$，中心分别为 $x_1,\ldots,x_N$。在\eqref{align:ch02-cube-image-estimate}中，以 $Q_j$ 代替 $Q$，并取 $T=D_{x_j}G$，得到
\begin{align*}
m(G(Q))&\le\sum_1^N m(G(Q_j))\\
&\le\sum_1^N|\det D_{x_j}G|\left(\sup_{y\in Q_j}\|(D_{x_j}G)^{-1}D_yG\|\right)^n m(Q_j)\\
&\le(1+\epsilon)\sum_1^N|\det D_{x_j}G|m(Q_j).
\end{align*}
最后这个和是 $\sum_1^N|\det D_{x_j}G|\chi_{Q_j}(x)$ 的积分；由于 $D_xG$ 连续，当 $\delta\to0$ 时，该函数在 $Q$ 上一致趋于 $|\det D_xG|$。因此，令 $\delta\to0$ 且 $\epsilon\to0$，得到
\begin{align*}
m(G(Q))\le\int_Q|\det D_xG|\,\mathrm{d}x.
\end{align*}
我们断言，以 $\Omega$ 中任意 Borel 集代替 $Q$，这一估计仍然成立。事实上，若 $U\subseteq\Omega$ 为开集，由\cref{lemma:ch02-open-cubes}，可写 $U=\bigcup_1^\infty Q_j$，其中 $Q_j$ 为内部两两不交的立方体。由于立方体的边界为 Lebesgue 零测集，
\begin{align*}
m(G(U))\le\sum_1^\infty m(G(Q_j))
\le\sum_1^\infty\int_{Q_j}|\det D_xG|\,\mathrm{d}x
=\int_U|\det D_xG|\,\mathrm{d}x.
\end{align*}
其次，令 $W_K=\Omega\cap\{x:|x|<K\text{ 且 }|\det D_xG|<K\}$。若 $E$ 为 $W_K$ 的 Borel 子集，由\cref{theorem:ch02-nd-regularity}，存在开集的递减序列 $U_j\subseteq W_{K+1}$，使 $E\subseteq\bigcup_1^\infty U_j$ 且 $m(\bigcap_1^\infty U_j\setminus E)=0$。由前面的估计和控制收敛定理，
\begin{align*}
m(G(E))&\le m\left(G\left(\bigcap_1^\infty U_j\right)\right)=\lim m(G(U_j))\\
&\le\lim\int_{U_j}|\det D_xG|\,\mathrm{d}x=\int_E|\det D_xG|\,\mathrm{d}x.
\end{align*}
最后，若 $E$ 为 $\Omega$ 的任意 Borel 子集，对 $E\cap W_K$ 应用这一论证，令 $K\to\infty$，再由单调收敛定理，得到 $m(G(E))\le\int_E|\det D_xG|\,\mathrm{d}x$。

因此，若 $f=\sum a_j\chi_{A_j}$ 为 $G(\Omega)$ 上的非负简单函数，则
\begin{align*}
\int_{G(\Omega)}f(x)\,\mathrm{d}x
&=\sum a_jm(A_j)\le\sum a_j\int_{G^{-1}(A_j)}|\det D_xG|\,\mathrm{d}x\\
&=\int_\Omega f\circ G(x)|\det D_xG|\,\mathrm{d}x.
\end{align*}
\cref{theorem:ch02-simple-approximation}与单调收敛定理随即表明，对任意非负可测函数 $f$，有
\begin{align*}
\int_{G(\Omega)}f(x)\,\mathrm{d}x\le\int_\Omega f\circ G(x)|\det D_xG|\,\mathrm{d}x.
\end{align*}
但以 $G^{-1}$ 代替 $G$、以 $f\circ G$ 代替 $f$，同样的推理仍适用，故
\begin{align*}
&\int_\Omega f\circ G(x)|\det D_xG|\,\mathrm{d}x\\
&\quad\le\int_{G(\Omega)}f\circ G\circ G^{-1}(x)|\det D_{G^{-1}(x)}G||\det D_xG|\,\mathrm{d}x
=\int_{G(\Omega)}f(x)\,\mathrm{d}x.
\end{align*}
这就证明了 $f\ge0$ 时的\subcref{theorem:ch02-differentiable-change}{item:ch02-differentiable-change-integrals}，$f\in L^1$ 的情形立即随之成立。而\subcref{theorem:ch02-differentiable-change}{item:ch02-differentiable-change-sets}正是\subcref{theorem:ch02-differentiable-change}{item:ch02-differentiable-change-integrals}中取 $f=\chi_{G(E)}$ 的特例，证明完成。
\end{proof}
\subsection*{习题}
% SOURCE p=91; printed=76; object=ch02-exercise-nd-density-proof; number=53; region=bottom
\begin{exercise}\label{exercise:ch02-nd-density-proof}
补全\cref{theorem:ch02-nd-lone-density}证明的细节。
\end{exercise}
% SOURCE p=92; printed=77; object=ch02-exercise-singular-linear-change; number=54; region=top
\begin{exercise}\label{exercise:ch02-singular-linear-change}
若 $T$ 不可逆，\cref{theorem:ch02-linear-change}中有多少结论仍然成立？
\end{exercise}
% SOURCE p=92; printed=77; object=ch02-exercise-iterated-integral-examples; number=55; region=upper
\begin{exercise}\label{exercise:ch02-iterated-integral-examples}
设 $E=[0,1]\times[0,1]$。对下列函数 $f$，研究 $\int_E f\,\mathrm{d}m^2$、$\int_0^1\int_0^1 f(x,y)\,\mathrm{d}x\,\mathrm{d}y$ 以及 $\int_0^1\int_0^1 f(x,y)\,\mathrm{d}y\,\mathrm{d}x$ 的存在性及相等关系。
\begin{enumerate}
\item\label{item:ch02-iterated-rational-example} $f(x,y)=(x^2-y^2)(x^2+y^2)^{-2}$。
\item\label{item:ch02-iterated-power-example} $f(x,y)=(1-xy)^{-a}$，$a>0$。
\item\label{item:ch02-iterated-strip-example} 当 $0<y<|x-\frac12|$ 时，$f(x,y)=(x-\frac12)^{-3}$；其余情形 $f(x,y)=0$。
\end{enumerate}
\end{exercise}
% SOURCE p=92; printed=77; object=ch02-exercise-integral-tail-average; number=56; region=upper
\begin{exercise}\label{exercise:ch02-integral-tail-average}
若 $f$ 在 $(0,a)$ 上 Lebesgue 可积，且 $g(x)=\int_x^a t^{-1}f(t)\,\mathrm{d}t$，则 $g$ 在 $(0,a)$ 上可积，且 $\int_0^a g(x)\,\mathrm{d}x=\int_0^a f(x)\,\mathrm{d}x$。
\end{exercise}
% SOURCE p=92; printed=77; object=ch02-exercise-laplace-sine-fubini; number=57; region=upper-middle
\begin{exercise}\label{exercise:ch02-laplace-sine-fubini}
通过关于 $x$ 和 $y$ 积分 $e^{-sxy}\sin x$，证明当 $s>0$ 时，$\int_0^\infty e^{-sx}x^{-1}\sin x\,\mathrm{d}x=\arctan(s^{-1})$。（回顾 $\tan(\frac\pi2-\theta)=(\tan\theta)^{-1}$ 可能有帮助。比较\subcref{exercise:ch02-termwise-integrals}{item:ch02-laplace-sine}。）
\end{exercise}
% SOURCE p=92; printed=77; object=ch02-exercise-laplace-sine-square; number=58; region=middle
\begin{exercise}\label{exercise:ch02-laplace-sine-square}
通过关于 $x$ 和 $y$ 积分 $e^{-sx}\sin2xy$，证明当 $s>0$ 时，$\int e^{-sx}x^{-1}\sin^2x\,\mathrm{d}x=\frac14\log(1+4s^{-2})$。
\end{exercise}
% SOURCE p=92; printed=77; object=ch02-exercise-dirichlet-integral; number=59; region=middle
\begin{exercise}\label{exercise:ch02-dirichlet-integral}
令 $f(x)=x^{-1}\sin x$。
\begin{enumerate}
\item\label{item:ch02-dirichlet-not-absolute} 证明 $\int_0^\infty|f(x)|\,\mathrm{d}x=\infty$。
\item\label{item:ch02-dirichlet-value} 通过关于 $x$ 和 $y$ 积分 $e^{-xy}\sin x$，证明 $\lim_{b\to\infty}\int_0^b f(x)\,\mathrm{d}x=\frac12\pi$。（鉴于\subcref{exercise:ch02-dirichlet-integral}{item:ch02-dirichlet-not-absolute}，令 $b\to\infty$ 取极限时需要小心。）
\end{enumerate}
\end{exercise}
% SOURCE p=92; printed=77; object=ch02-exercise-beta-integral; number=60; region=middle-lower
\begin{exercise}\label{exercise:ch02-beta-integral}
对 $x,y>0$，$\Gamma(x)\Gamma(y)/\Gamma(x+y)=\int_0^1t^{x-1}(1-t)^{y-1}\,\mathrm{d}t$。（回顾 $\Gamma$ 是在\cref{section:ch02-complex-integration}中定义的。将 $\Gamma(x)\Gamma(y)$ 写成二重积分，并以指数函数的自变量作为新的积分变量。）
\end{exercise}
% SOURCE p=92; printed=77; object=ch02-exercise-fractional-integral; number=61; region=lower
\begin{exercise}\label{exercise:ch02-fractional-integral}
若 $f$ 在 $[0,\infty)$ 上连续，
\begin{definition*}[分数阶积分]
对 $\alpha>0$ 和 $x\ge0$，令
\begin{align*}
I_\alpha f(x)=\frac1{\Gamma(\alpha)}\int_0^x(x-t)^{\alpha-1}f(t)\,\mathrm{d}t.
\end{align*}
$I_\alpha f$ 称为 $f$ 的 \textbf{$\alpha$ 阶分数积分}。
\end{definition*}
\begin{enumerate}
\item\label{item:ch02-fractional-integral-semigroup} 对所有 $\alpha,\beta>0$，$I_{\alpha+\beta}f=I_\alpha(I_\beta f)$。（使用\cref{exercise:ch02-beta-integral}。）
\item\label{item:ch02-integer-integral-antiderivative} 若 $n\in\mathbb{N}$，则 $I_nf$ 为 $f$ 的 $n$ 阶原函数。
\end{enumerate}
\end{exercise}
\end{document}
